\section{Model}
\label{sec:axm-model}

Time is continuous, and I suppress time subscripts except where they are needed
for clarity. There is one final good, a representative household, competitive
final-good firms, and an integrated AI developer. Two productivity indexes enter
the model. Labor productivity \(A\) grows exogenously. AI capability \(B\) is
produced endogenously by the developer. The developer monopolistically supplies
AI production services and invests in better AI. The initial economy already
uses AI production services; the model does not describe their first commercial
adoption.

I first present a benchmark in which AI research services are the only input
used to improve AI capability. This benchmark represents an economy in which
the replacement of human researchers in frontier AI development has already
occurred. I then introduce human research as an extension. This order separates
the mechanism needed for the main results from the additional mechanism that
affects the timing of the transition.

\tikzset{
 axm-model-input/.style={draw, rounded corners=2pt, minimum width=2.7cm,
  minimum height=1.0cm, align=center, fill=gray!8},
 axm-model-state/.style={draw, rounded corners=2pt, minimum width=2.7cm,
  minimum height=1.0cm, align=center, fill=blue!8},
 axm-model-technology/.style={draw, rounded corners=2pt, minimum width=3.5cm,
  minimum height=1.0cm, align=center, fill=orange!9},
 axm-model-problem/.style={draw, rounded corners=2pt, minimum height=1.2cm,
  align=center, fill=gray!6, text width=10.8cm},
 axm-model-flow/.style={->, semithick},
 axm-model-dynamics/.style={->, semithick, dashed}
}

\subsection{Population, labor endowment, and labor productivity}

Population grows exogenously at the constant rate \(n>0\):
\begin{equation}
 \dot N=nN.
 \label{eq:axm-population}
\end{equation}
Each person supplies one unit of homogeneous labor inelastically, so the
household supplies \(N\) units at the common wage \(w\). In the benchmark, all
labor is used in final production. The extension later allows the AI developer
to hire part of the labor endowment for research.

The variable \(A\) measures labor-augmenting productivity in final production.
One unit of production labor supplies \(A\) efficiency units, so effective
production labor is \(AL\). Labor productivity follows the exogenous law
\begin{equation}
 \dot A=\gamma_A A,
 \qquad \gamma_A\geq0,
 \label{eq:axm-labor-productivity}
\end{equation}
with \(A(0)>0\). The case \(\gamma_A=0\) nests the specification without
conventional labor-productivity growth. For any positive variable \(V\), write
\(g_V\equiv\dot V/V\); hence \(g_A=\gamma_A\).

\subsection{Household}

The representative household discounts utility at rate \(\rho>n\). It takes the
paths of the net interest rate \(r\), the wage \(w\), distributed developer
profit \(\Pi\), and population \(N\) as given. Let \(K\) denote its holdings of
physical capital. Given \(K(0)>0\), it chooses \(C\) and \(K\) to maximize
\[
 \int_0^\infty e^{-\rho t}N_t\log(C_t/N_t)\,\dd t.
\]
It owns physical capital and the AI developer and receives labor income, the net
return on capital, and distributed profit. Its flow budget is
\begin{equation}
 \dot K=rK+wN+\Pi-C.
 \label{eq:axm-household-budget}
\end{equation}
The variable \(r\) is the net real interest rate: depreciation is already
deducted from the rental return to capital. At this stage, \(r\) is a price the
household takes as given; no production-side condition has yet been imposed. The
household's interior Euler equation in aggregate variables is therefore
\begin{equation}
 \frac{\dot C}{C}=n+r-\rho.
 \label{eq:axm-euler}
\end{equation}
The infinite-horizon terminal restriction maintained for the household is
\begin{equation}
 \lim_{t\to\infty}e^{-\rho t}\frac{N_tK_t}{C_t}=0.
 \label{eq:axm-household-tvc}
\end{equation}

\subsection{Final-good firms}

Competitive final-good firms take \(r,w,p_X\) as given and demand capital,
production labor, and AI production services, where \(p_X\) is their price. To
preserve the notation used in the analysis and numerical code, I do not give
\(K\) and \(X\) agent-specific superscripts. In this subsection they denote
capital and AI production services demanded by the final sector; in the
household problem \(K\) denotes capital holdings, and in the developer's problem
\(X\) denotes sales. Equality of the corresponding quantities is imposed only
in equilibrium.

Let \(0<\alpha<1\) be the physical-capital elasticity, let \(\delta\geq0\) be
the depreciation rate, and let \(\sigma_{XL}>0\) be the elasticity of
substitution between effective production labor and AI production services. The
CES weights satisfy \(\omega_L,\omega_X\in(0,1)\) and
\(\omega_L+\omega_X=1\). The production technology is
\begin{align}
 Y&=K^\alpha Z^{1-\alpha},\\
 Z&=\left[\omega_L (AL)^{\frac{\sigma_{XL}-1}{\sigma_{XL}}}
 +\omega_X X^{\frac{\sigma_{XL}-1}{\sigma_{XL}}}\right]^{
 \frac{\sigma_{XL}}{\sigma_{XL}-1}},
 \qquad \omega_L+\omega_X=1.
 \label{eq:axm-final-production}
\end{align}
At the unit-elasticity limit, \(\sigma_{XL}\to1\), the CES becomes
\(Z=(AL)^{\omega_L}X^{\omega_X}\). Define the contribution of AI production
services to the CES quantity index as
\begin{equation}
 s_X=\frac{\omega_X X^{(\sigma_{XL}-1)/\sigma_{XL}}}
 {\omega_L (AL)^{(\sigma_{XL}-1)/\sigma_{XL}}
 +\omega_X X^{(\sigma_{XL}-1)/\sigma_{XL}}}.
 \label{eq:axm-sx}
\end{equation}
Thus \(s_X\) is the elasticity of \(Z\) with respect to \(X\). The competitive
sector solves
\begin{equation}
 \max_{K,L,X\geq0}\left\{K^\alpha Z(AL,X)^{1-\alpha}
 -(r+\delta)K-wL-p_XX\right\}.
 \label{eq:axm-final-firm-problem}
\end{equation}
The interior first-order conditions are
\begin{equation}
 r=\alpha\frac{Y}{K}-\delta,\qquad
 w=(1-\alpha)(1-s_X)\frac{Y}{L},\qquad
 p_X=(1-\alpha)s_X\frac{Y}{X}.
 \label{eq:axm-factor-prices}
\end{equation}
Labor productivity \(A\) affects the wage through \(Y\) and \(s_X\). It does
not appear as an additional factor in the wage equation because \(AL\) is
linear in \(L\), so the output elasticity with respect to \(L\) equals the
elasticity with respect to effective labor \(AL\).

To ask how an expansion of AI production services affects production labor,
holding \(K\) fixed is not enough: either the wage or labor must also be held
fixed. I therefore state both sides of the same comparative static. The
conditional labor-demand response holds \(K\), \(A\), and \(w\) fixed; the dual
wage response holds \(K\), \(A\), and \(L\) fixed. The distinction between the
wage and labor-share responses follows the general result in
\citet[Theorem~1]{autorkausik2026}. For this nested Cobb--Douglas--CES
technology, \citet[Proposition~2.2, Appendix~A.4, Lemma~4.1, and
Proposition~4.2]{yango2026} derive the corresponding labor-demand sign and the
two elasticity thresholds.

Define production labor's share of final output as
\begin{equation}
 \ell^Y\equiv\frac{wL}{Y}=(1-\alpha)(1-s_X).
 \label{eq:axm-production-labor-share}
\end{equation}

\Needspace{24\baselineskip}
\begin{proposition}[AI services and production labor]
\label{prop:axm-ai-labor-comparative-static}
For an interior final-good firm, an increase in \(X\), holding \(K\), \(A\),
and \(L\) fixed, satisfies
\begin{align}
 \left.\frac{\partial\ln Y}{\partial\ln X}\right|_{K,A,L}
 &=(1-\alpha)s_X>0,
 \label{eq:axm-output-x-elasticity}\\
 \left.\frac{\partial\ln w}{\partial\ln X}\right|_{K,A,L}
 &=\left(\frac{1}{\sigma_{XL}}-\alpha\right)s_X,
 \label{eq:axm-wage-x-elasticity}\\
 \left.\frac{\partial\ln \ell^Y}{\partial\ln X}\right|_{K,A,L}
 &=-\left(1-\frac{1}{\sigma_{XL}}\right)s_X.
 \label{eq:axm-labor-share-x-elasticity}
\end{align}
Equivalently, holding \(K\), \(A\), and \(w\) fixed, conditional
production-labor demand satisfies
\begin{equation}
 \left.\frac{\partial\ln L}{\partial\ln X}\right|_{K,A,w}
 =\frac{s_X\left(\sigma_{XL}^{-1}-\alpha\right)}
 {\alpha(1-s_X)+s_X/\sigma_{XL}}.
 \label{eq:axm-labor-demand-x-elasticity}
\end{equation}
Thus an expansion of AI production services raises conditional labor demand---or,
at fixed labor, raises the wage---if and only if
\(\sigma_{XL}<1/\alpha\). It lowers conditional labor demand and the wage if and
only if \(\sigma_{XL}>1/\alpha\). By contrast, production labor's share falls if
and only if \(\sigma_{XL}>1\). Since \(0<\alpha<1\), the interval
\begin{equation}
 1<\sigma_{XL}<\frac{1}{\alpha}
 \label{eq:axm-wage-share-decoupling-region}
\end{equation}
is nonempty: within it, more AI production services raise labor demand at a
given wage, and raise the wage at a given employment level, even as labor
receives a smaller share of output.
\end{proposition}

\begin{proof}
The CES share in \eqref{eq:axm-sx} implies
\[
 \left.\frac{\partial\ln(1-s_X)}{\partial\ln X}\right|_{A,L}
 =-\left(1-\frac{1}{\sigma_{XL}}\right)s_X.
\]
The output elasticity in \eqref{eq:axm-output-x-elasticity} follows directly
from \eqref{eq:axm-final-production}. Combining these two elasticities with the
wage equation in \eqref{eq:axm-factor-prices} gives
\eqref{eq:axm-wage-x-elasticity}; combining the share equation above with
\eqref{eq:axm-production-labor-share} gives
\eqref{eq:axm-labor-share-x-elasticity}. Finally,
\[
 \left.\frac{\partial\ln w}{\partial\ln L}\right|_{K,A,X}
 =-\left[\alpha(1-s_X)+\frac{s_X}{\sigma_{XL}}\right]<0.
\]
Implicit differentiation of the labor-demand condition at fixed \(w\) gives
\eqref{eq:axm-labor-demand-x-elasticity}. Its denominator is strictly positive,
so its sign is the sign of \(\sigma_{XL}^{-1}-\alpha\). The remaining sign
statements follow from the displayed elasticities.
\end{proof}

The final-good firm's AI-demand condition implies
\(p_XX/Y=(1-\alpha)s_X\). After market clearing, this is the developer's gross
service revenue as a share of final output. Constant returns and the three
first-order conditions imply zero profit:
\begin{equation}
 Y=(r+\delta)K+wL+p_XX.
 \label{eq:axm-final-zero-profit}
\end{equation}
Solving the AI-demand condition for its price yields the inverse demand schedule
\(p_X(X;K,A,L)\). Define its inverse absolute elasticity as
\begin{equation}
 \varepsilon_X^{-1}
 \equiv-\left.\frac{\partial\ln p_X(X;K,A,L)}{\partial\ln X}\right|_{K,A,L}
 =\frac{1-s_X}{\sigma_{XL}}+\alpha s_X.
 \label{eq:axm-inverse-elasticity}
\end{equation}

\begin{remark}[Tasks, jobs, and labor services]
The input \(L\) is a quantity of homogeneous labor services, not a number of
tasks, occupations, or jobs. The benchmark imposes full employment in final
production. The extension allows workers to move between production and
research. Neither version describes unemployment or the number of tasks
displaced by AI.
\end{remark}

\subsection{Integrated AI developer: automated-research benchmark}

\subsubsection{AI production and research services}

The model distinguishes two uses of compute. \(U\) is inference compute: raw
compute used to operate existing AI systems and supply services to final-good
firms. \(M\) is AI R\&D compute---called research compute below---used to improve
AI capability. Both are flows rather than stocks of chips, and both have the
same constant unit resource cost. I choose units so that one unit of raw compute
costs one unit of final output.

The stock \(B\) indexes current AI capability: what one unit of compute can
accomplish. Inference compute produces the flow of AI production services
\begin{equation}
 X=BU.
 \label{eq:axm-ai-service}
\end{equation}
Research compute analogously produces the flow of AI research services \(BM\).
The paid inputs are \(U\) and \(M\), whereas \(X\) and \(BM\) are service flows.
This product of nonrival capability and rival compute is the one-capability
version of the capability--compute accounting in \citet{davidsonetal2026}.
Producing one unit of either service flow uses \(1/B\) units of final output.

Figure~\ref{fig:axm-current-production-flow} collects the current-production
technology. Gray boxes denote current inputs, blue boxes denote stocks, and
orange boxes denote transformations or outputs. The diagram is a schematic,
not a numerical result.

\begin{figure}[H]
 \centering
 \resizebox{0.90\textwidth}{!}{%
 \begin{tikzpicture}[>=Latex,font=\small]
  \node[axm-model-state]      (B) at (-8.0,-0.2) {\(B\)\\AI capability};
  \node[axm-model-input]      (U) at (-8.0,-1.6) {\(U\)\\inference compute};
  \node[axm-model-technology] (X) at (-4.5,-0.9)
   {\(X=BU\)\\AI production services};
  \node[axm-model-state]      (A) at (-8.0, 2.7) {\(A\)\\labor productivity};
  \node[axm-model-input]      (L) at (-8.0, 1.3)
   {\(L\)\\production labor};
  \node[axm-model-technology] (AL) at (-4.5, 2.0)
   {\(AL\)\\effective labor};
  \node[axm-model-technology] (Z) at ( 0.0, 0.8)
   {\(Z=\operatorname{CES}(AL,X)\)\\elasticity \(\sigma_{XL}\)};
  \node[axm-model-state]      (K) at ( 0.0, 3.7)
   {\(K\)\\physical capital};
  \node[axm-model-technology] (Y) at ( 4.0, 2.5)
   {\(Y=K^\alpha Z^{1-\alpha}\)\\final output};

  \draw[axm-model-flow] (B) -- (X);
  \draw[axm-model-flow] (U) -- (X);
  \draw[axm-model-flow] (A) -- (AL);
  \draw[axm-model-flow] (L) -- (AL);
  \draw[axm-model-flow] (AL) -- (Z);
  \draw[axm-model-flow] (X) -- (Z);
  \draw[axm-model-flow] (Z) -- (Y);
  \draw[axm-model-flow] (K) -- (Y);
 \end{tikzpicture}%
 }
 \caption{Current production. Labor productivity and production labor produce
 effective labor \(AL\). AI capability and inference compute produce AI
 production services \(X\). Effective labor and AI production services form the
 CES composite \(Z\), which combines with physical capital to produce final
 output.}
 \label{fig:axm-current-production-flow}
\end{figure}

\subsubsection{Capability accumulation}

In the benchmark, AI research services are the only input used to improve
capability:
\begin{equation}
 \dot B=\chi(BM)^\eta,
 \qquad \chi>0,\quad 0<\eta<1.
 \label{eq:axm-capability-law-benchmark}
\end{equation}
The scalar \(\chi\) shifts research productivity. The exponent \(\eta\) governs
returns to current research scale: for a fixed \(B\), multiplying \(M\) by
\(\lambda\) multiplies \(\dot B\) by \(\lambda^\eta\). Capability enters only
through \(BM\), so the benchmark has no separate \(B^\phi\) term. Better
capability makes each unit of research compute more productive, while \(\eta\)
governs decreasing returns to the current flow of research compute.

\begin{figure}[H]
 \centering
 \resizebox{0.78\textwidth}{!}{%
 \begin{tikzpicture}[>=Latex,font=\small]
  \node[axm-model-state]      (B)  at (-5.7, 1.2)
   {\(B\)\\AI capability};
  \node[axm-model-input]      (M)  at (-5.7,-1.2)
   {\(M\)\\research compute};
  \node[axm-model-technology] (BM) at (-1.7, 0.0)
   {\(BM\)\\AI research services};
  \node[axm-model-technology] (Bd) at ( 3.2, 0.0)
   {\(\dot B=\chi(BM)^\eta\)\\new capability};

  \draw[axm-model-flow] (B) -- (BM);
  \draw[axm-model-flow] (M) -- (BM);
  \draw[axm-model-flow] (BM) -- (Bd);
  \draw[axm-model-dynamics] (Bd.east) -- (5.2,0.0) -- (5.2,2.8) --
   node[above] {capability accumulation} (-5.7,2.8) -- (B.north);
 \end{tikzpicture}%
 }
 \caption{AI research in the benchmark. Capability and research compute produce
 AI research services \(BM\). Those services improve capability, and the new
 capability stock raises the productivity of subsequent research compute.}
 \label{fig:axm-research-flow}
\end{figure}

\subsubsection{Dynamic problem}

The developer sells \(X\), pays \(X/B\) for inference compute, uses \(M\) for
research, and owns \(B\). It faces the inverse demand schedule derived from
final-good firms' optimal demand. Taking the paths of \(r,K,A,\) and \(L\) as
given, it solves
\begin{equation}
 \max_{\{X,M\}}
 \int_0^\infty
 \exp\left\{-\int_0^t r_s\dd s\right\}
 \left[p_X(X_t;K_t,A_t,L_t)X_t-\frac{X_t}{B_t}-M_t\right]\dd t,
 \label{eq:axm-developer-problem-benchmark}
\end{equation}
subject to \eqref{eq:axm-capability-law-benchmark} and the initial capability
stock. The discount rate is the market return available to the household that
owns the developer. Equilibrium later requires consistency between this return,
household saving, and final-good firms' capital demand. The developer's
instantaneous distributed profit is
\begin{equation}
 \Pi=p_X(X;K,A,L)X-\frac{X}{B}-M=p_XX-U-M.
 \label{eq:axm-developer-profit-benchmark}
\end{equation}
This net payout can be negative. The model imposes no separate cash-flow or
borrowing constraint, so a negative \(\Pi\) represents an equity injection from
the household that owns the developer.

Sales \(X\) do not enter the capability law or another intertemporal constraint,
and there is no firm-level capacity constraint linking \(U\) and \(M\). The
developer therefore chooses \(X\) pointwise, while \(M\) affects future
capability. Conditional on current \(B,K,A,\) and \(L\), optimized operating
profit before research expenditure is
\begin{equation}
 \pi(B;K,A,L)=\max_{X\geq0}
 \left\{p_X(X;K,A,L)X-\frac{X}{B}\right\}.
 \label{eq:axm-service-supply}
\end{equation}
On an interior monopoly branch, the first-order condition is
\begin{equation}
 p_X(1-\varepsilon_X^{-1})=\frac{1}{B}.
 \label{eq:axm-monopoly-foc}
\end{equation}
Writing \(e_X=\varepsilon_X^{-1}\), the strict second-order condition is
\begin{equation}
 e_X(1-e_X)
 +(\alpha-\sigma_{XL}^{-1})(1-\sigma_{XL}^{-1})s_X(1-s_X)>0.
 \label{eq:axm-monopoly-soc}
\end{equation}

Whenever the pointwise maximum is attained, the primitive problem is equivalent
to
\begin{equation}
 \max_{\{M\}}
 \int_0^\infty
 \exp\left\{-\int_0^t r_s\dd s\right\}
 \left[\pi(B_t;K_t,A_t,L_t)-M_t\right]\dd t,
 \label{eq:axm-developer-problem-reduced-benchmark}
\end{equation}
subject to the same capability law and initial stock. The maximization over
\(X\) is therefore the pointwise inner maximization of one dynamic problem, not
a second decision problem.

\begin{figure}[H]
 \centering
 \begin{tikzpicture}[>=Latex,font=\small,node distance=1.25cm]
  \node[axm-model-problem] (full) {
   \textbf{Primitive dynamic problem, \eqref{eq:axm-developer-problem-benchmark}}\\
   choose \(X,M\); state \(B\); flow payoff
   \(p_X(X;K,A,L)X-X/B-M\)};
  \node[axm-model-problem, below=of full, fill=orange!9] (inner) {
   \textbf{Pointwise inner choice, \eqref{eq:axm-service-supply}}\\
   \(\displaystyle \pi(B;K,A,L)=\max_{X\geq0}
   \{p_X(X;K,A,L)X-X/B\}\)};
  \node[axm-model-problem, below=of inner, fill=blue!8] (reduced) {
   \textbf{Equivalent reduced dynamic problem,
   \eqref{eq:axm-developer-problem-reduced-benchmark}}\\
   choose \(M\); same state law and initial \(B\); flow payoff
   \(\pi(B;K,A,L)-M\)};

  \draw[axm-model-flow] (full) -- node[right,align=left]
   {\(X\) does not enter \(\dot B\) or another\\intertemporal constraint}
   (inner);
  \draw[axm-model-flow] (inner) -- node[right]
   {substitute optimized operating profit} (reduced);
 \end{tikzpicture}
 \caption{Decomposition of the developer's benchmark problem. The maximization
 over \(X\) is an inner, within-date choice. Substituting optimized operating
 profit leaves the dynamic choice of research compute \(M\).}
 \label{fig:axm-developer-decomposition}
\end{figure}

\subsubsection{Dynamic optimality}

Let \(q\) be the current-value shadow price of AI capability. The interior
first-order condition for research compute is
\begin{equation}
 q\chi\eta B(BM)^{\eta-1}=1.
 \label{eq:axm-research-foc-benchmark}
\end{equation}
Equivalently,
\begin{equation}
 M=\left(q\chi\eta B^\eta\right)^{1/(1-\eta)}.
 \label{eq:axm-research-demand-benchmark}
\end{equation}
This expression determines current research compute after \(B\) and \(q\) are
known. It is derived from the developer's problem and is not an additional
resource constraint.

The envelope theorem applied to \eqref{eq:axm-service-supply} gives
\(\pi_B=X/B^2\). Holding \(M\) fixed, the capability law gives
\(\partial\dot B/\partial B=\eta\dot B/B\). The current-value costate equation is
therefore
\begin{equation}
 \frac{\dot q}{q}=r-\frac{X}{qB^2}-\eta\frac{\dot B}{B}.
 \label{eq:axm-costate-benchmark}
\end{equation}
The infinite-horizon terminal restriction maintained for the developer is
\begin{equation}
 \lim_{t\to\infty}
 \exp\left\{-\int_0^t r_s\dd s\right\}q_tB_t=0.
 \label{eq:axm-developer-tvc}
\end{equation}

\subsection{Benchmark equilibrium}

Capital-market clearing requires the household's capital holdings to equal the
capital rented by final-good firms. Market clearing for AI production services
requires the quantity purchased by final-good firms to equal the quantity sold
by the developer. Labor-market clearing in the benchmark imposes
\begin{equation}
 N=L.
 \label{eq:axm-labor-market-benchmark}
\end{equation}

Substituting developer profit \eqref{eq:axm-developer-profit-benchmark}, labor-
market clearing, and market clearing for AI production services into the
household budget gives
\[
 \dot K=rK+wL+p_XX-U-M-C.
\]
Using the final sector's zero-profit identity yields the aggregate resource
constraint
\begin{equation}
 C+\dot K+\delta K+U+M=Y.
 \label{eq:axm-resource}
\end{equation}
The corresponding accounting decomposition is
\begin{align}
 Y&=(r+\delta)K+wL+\Pi+U+M,\notag\\
 1&=\frac{(r+\delta)K}{Y}+\frac{wL}{Y}
 +\frac{\Pi}{Y}+\frac{U}{Y}+\frac{M}{Y}.
 \label{eq:axm-consolidated-distribution-benchmark}
\end{align}
This identity separates the gross capital payment, labor income, distributed
developer profit, inference compute, and research compute. It is an accounting
consequence of equilibrium, not an additional restriction.

Only after market clearing can conditions from different agents be combined.
The household Euler equation and the final-good firm's capital demand imply
\begin{equation}
 \frac{\dot C}{C}=n+\alpha\frac{Y}{K}-\delta-\rho.
 \label{eq:axm-euler-reduced}
\end{equation}
Substituting the same capital-price condition into the developer's costate
equation gives
\begin{equation}
 \frac{\dot q}{q}=\alpha\frac{Y}{K}-\delta-
 \frac{X}{qB^2}-\eta\frac{\dot B}{B}.
 \label{eq:axm-costate-reduced-benchmark}
\end{equation}

The benchmark contains a technological loop,
\(B\to BM\to\dot B\to B\), and an economic loop,
\(B\to X\to Y\to(K,M)\to BM\to\dot B\to B\). Three local elasticities
summarize the direct technological links:
\begin{equation}
 \left.\frac{\partial\ln X}{\partial\ln B}\right|_U=1,\qquad
 \left.\frac{\partial\ln Y}{\partial\ln X}\right|_{K,A,L}
 =(1-\alpha)s_X,\qquad
 \frac{\partial\ln\dot B}{\partial\ln(BM)}=\eta.
 \label{eq:axm-local-feedback-links-benchmark}
\end{equation}
The strength of the full economic loop is not a primitive constant because
saving, compute allocation, and prices are determined in equilibrium.

\begin{definition}[Admissible path]
An admissible path has the measurable nonnegative controls defined by the
version of the model under consideration, absolutely continuous nonnegative
stocks \(K\) and \(B\), the exogenous paths of \(N\) and \(A\), and the stated
initial values. It satisfies the relevant flow budget and state equations, and
its improper objective integrals converge to finite real values.
\end{definition}

\begin{definition}[Interior benchmark equilibrium]
Given \(K(0),A(0),B(0),N(0)>0\), an interior benchmark equilibrium is an
admissible path with \(C,K,A,B,N,L,U,M,X,Z,Y>0\) and \(w,p_X,q>0\). The net
interest rate \(r\) and developer payout \(\Pi\) may have either sign. The
household, final-good firms, and developer solve their respective problems, and
labor, capital, AI production services, and goods markets clear.
\end{definition}

I call an equilibrium regular when the associated maximum-principle
multipliers are normal, the relevant paths are differentiable wherever the
displayed derivatives are used, and the pointwise maxima are attained. Under
these conditions, the dated first-order and costate equations apply almost
everywhere.

\subsubsection{Reduced benchmark system}

Given \(K,A,B,N,q\), the algebraic subsystem consists of \(L=N\), final
production and factor prices, the service identity \(X=BU\), inverse demand and
the monopoly first-order condition, and the research-compute condition
\eqref{eq:axm-research-demand-benchmark}. It determines
\(L,U,M,X,Z,Y,r,w,p_X\), and \(s_X\) on the selected monopoly branch. Combining
this algebraic block with the state and costate equations gives
\begin{subequations}
\label{eq:axm-canonical-system-benchmark}
\begin{align}
 \frac{\dot N}{N}&=n,\\
 \frac{\dot A}{A}&=\gamma_A,\\
 \frac{\dot K}{K}&=\frac{Y-C-U-M}{K}-\delta,\\
 \frac{\dot B}{B}&=\frac{\chi(BM)^\eta}{B},\\
 \frac{\dot C}{C}&=n+\alpha\frac{Y}{K}-\delta-\rho,\\
 \frac{\dot q}{q}
 &=\alpha\frac{Y}{K}-\delta-\frac{X}{qB^2}
 -\eta\frac{\dot B}{B}.
\end{align}
\end{subequations}
The algebraic subsystem may admit multiple roots. I follow a branch that
satisfies the monopoly second-order condition and make no claim of global
uniqueness.

After capital-market clearing, \(K\) and \(B\) are predetermined, the paths of
\(A\) and \(N\) are exogenous, and \(C\) and \(q\) are jump variables. The
initial values of \(K,A,B,N\), market clearing, and the two terminal restrictions
close the system. In equilibrium the terminal conditions can be collected as
\begin{equation}
 \lim_{t\to\infty}e^{-\rho t}\frac{N_tK_t}{C_t}=0,\qquad
 \lim_{t\to\infty}
 \exp\left\{-\int_0^t\left(\alpha\frac{Y_s}{K_s}-\delta\right)\dd s\right\}
 q_tB_t=0.
 \label{eq:axm-tvc}
\end{equation}
A finite-horizon numerical calculation cannot itself verify either
infinite-horizon limit.

\subsection{Extension: human research in capability production}
\label{sec:axm-human-research-extension}

I now allow the developer to hire human research labor \(H\) at the same wage
paid by final-good firms. Labor-market clearing becomes
\begin{equation}
 N=L+H.
 \label{eq:axm-labor-market}
\end{equation}
The productivity index \(A\) continues to augment production labor \(L\), not
human research labor \(H\).

Let \(\sigma_{HM}>0\) be the elasticity of substitution between human research
and AI research services. The CES weights satisfy
\(\omega_H,\omega_M\in(0,1)\) and \(\omega_H+\omega_M=1\). For
\(\sigma_{HM}\neq1\), capability evolves according to
\begin{equation}
 \dot B=\chi
 \left[\omega_H H^{\varrho_{HM}}
 +\omega_M(BM)^{\varrho_{HM}}\right]^{\eta/\varrho_{HM}},
 \qquad
 \varrho_{HM}=\frac{\sigma_{HM}-1}{\sigma_{HM}}.
 \label{eq:axm-capability-law}
\end{equation}
Define the constant-returns effective-research index
\begin{equation}
 E\equiv\left[\omega_H H^{\varrho_{HM}}
 +\omega_M(BM)^{\varrho_{HM}}\right]^{1/\varrho_{HM}}.
 \label{eq:axm-effective-research}
\end{equation}
Then \(\dot B=\chi E^\eta\). At the unit-elasticity limit,
\begin{equation}
 \dot B=\chi H^{\eta\omega_H}(BM)^{\eta\omega_M},
 \qquad E=H^{\omega_H}(BM)^{\omega_M}.
 \label{eq:axm-capability-law-cd}
\end{equation}
The elasticity \(\sigma_{HM}\) governs substitution between the two research
inputs; \(\eta\) continues to govern returns to their joint scale. Setting
\(\omega_H=0\) and \(\omega_M=1\) recovers the benchmark
\eqref{eq:axm-capability-law-benchmark}; at that boundary
\(\sigma_{HM}\) is irrelevant.

The developer's primitive problem becomes
\begin{equation}
 \max_{\{X,H,M\}}
 \int_0^\infty
 \exp\left\{-\int_0^t r_s\dd s\right\}
 \left[p_X(X_t;K_t,A_t,L_t)X_t-\frac{X_t}{B_t}-w_tH_t-M_t\right]\dd t,
 \label{eq:axm-developer-problem}
\end{equation}
subject to \eqref{eq:axm-capability-law} and the initial capability stock. Its
instantaneous distributed profit is
\begin{equation}
 \Pi=p_XX-U-wH-M.
 \label{eq:axm-developer-profit}
\end{equation}
The pointwise choice of \(X\) is unchanged. Substituting optimized operating
profit gives the equivalent reduced problem
\begin{equation}
 \max_{\{H,M\}}
 \int_0^\infty
 \exp\left\{-\int_0^t r_s\dd s\right\}
 \left[\pi(B_t;K_t,A_t,L_t)-w_tH_t-M_t\right]\dd t,
 \label{eq:axm-developer-problem-reduced}
\end{equation}
subject to the same capability law and initial stock.

\subsubsection{Conditional research cost}

Conditional on current \(B\) and \(w\), the unit-expenditure problem for \(E\)
is
\begin{equation}
 P_E(B,w)=\min_{H,M\geq0}\left\{wH+M:E(H,BM)\geq1\right\}.
 \label{eq:axm-research-cost-problem}
\end{equation}
This is the dual representation of the research-input choice inside
\eqref{eq:axm-developer-problem-reduced}. For \(\sigma_{HM}\neq1\), its solution
is
\begin{equation}
 P_E=\left[\omega_H^{\sigma_{HM}}w^{1-\sigma_{HM}}
 +\omega_M^{\sigma_{HM}}\left(\frac{1}{B}\right)^{1-\sigma_{HM}}
 \right]^{1/(1-\sigma_{HM})}.
 \label{eq:axm-research-price}
\end{equation}
At the unit-elasticity limit,
\begin{equation}
 P_E=\left(\frac{w}{\omega_H}\right)^{\omega_H}
 \left(\frac{1}{B\omega_M}\right)^{\omega_M}.
 \label{eq:axm-research-price-cd}
\end{equation}
The minimum expenditure required to produce a target \(\dot B\) is
\begin{equation}
 P_E\left(\frac{\dot B}{\chi}\right)^{1/\eta}.
 \label{eq:axm-research-expenditure-function}
\end{equation}
Thus \(P_E\) is the unit cost of \(E\), not the unit cost of \(\dot B\) when
\(\eta<1\). Conditional input demands are
\begin{equation}
 H=\omega_H^{\sigma_{HM}}\left(\frac{P_E}{w}\right)^{\sigma_{HM}}E,
 \qquad
 BM=\omega_M^{\sigma_{HM}}(BP_E)^{\sigma_{HM}}E.
 \label{eq:axm-research-demands}
\end{equation}

The share of research expenditure allocated to AI research services is
\begin{equation}
 s_M=\frac{M}{wH+M}.
 \label{eq:axm-sm}
\end{equation}
Along an interior cost-minimizing allocation, \(s_M\) is also the elasticity of
\(E\) with respect to \(BM\). The elasticity of \(\dot B\) with respect to
\(BM\) is therefore \(\eta s_M\). Holding \(H\) and \(M\) fixed, \(\eta s_M\)
is also the local elasticity of \(\dot B\) with respect to \(B\). The variable
\(s_M\) is an expenditure share and a local feedback weight, not a physical
fraction of research tasks performed by AI.

\subsubsection{Dynamic optimality and equilibrium changes}
\label{sec:axm-dated-equilibrium-system}

The direct first-order conditions for the two research inputs are
\begin{equation}
 q\frac{\partial\dot B}{\partial H}=w,\qquad
 q\frac{\partial\dot B}{\partial M}=1.
 \label{eq:axm-research-focs}
\end{equation}
Equivalently, conditional cost minimization and the optimal research scale give
\begin{equation}
 q\chi\eta E^{\eta-1}=P_E.
 \label{eq:axm-research-scale-foc}
\end{equation}
The current-value costate equation becomes
\begin{equation}
 \frac{\dot q}{q}=r-\frac{X}{qB^2}
 -\eta s_M\frac{\dot B}{B}.
 \label{eq:axm-costate}
\end{equation}
Relative to the benchmark, human research changes the last term because only
the AI-research component \(BM\) is directly augmented by current capability.

The resource constraint \eqref{eq:axm-resource} is unchanged: human research
uses labor from the fixed labor endowment, while research compute \(M\) uses the
final good. The consolidated accounting identity now separates the two wage
bills:
\begin{align}
 Y&=(r+\delta)K+wL+wH+\Pi+U+M,\notag\\
 1&=\frac{(r+\delta)K}{Y}+\frac{wL}{Y}+\frac{wH}{Y}
 +\frac{\Pi}{Y}+\frac{U}{Y}+\frac{M}{Y}.
 \label{eq:axm-consolidated-distribution}
\end{align}
The benchmark identity follows by setting \(H=0\).

After imposing market clearing, the reduced costate equation is
\begin{equation}
 \frac{\dot q}{q}=\alpha\frac{Y}{K}-\delta-\frac{X}{qB^2}
 -\eta s_M\frac{\dot B}{B}.
 \label{eq:axm-costate-reduced}
\end{equation}
The extension replaces the benchmark's research elasticity \(\eta\) with the
endogenous feedback weight \(\eta s_M\):
\begin{equation}
 \left.\frac{\partial\ln X}{\partial\ln B}\right|_U=1,\qquad
 \left.\frac{\partial\ln Y}{\partial\ln X}\right|_{K,A,L}
 =(1-\alpha)s_X,\qquad
 \left.\frac{\partial\ln\dot B}{\partial\ln(BM)}\right|_H=\eta s_M.
 \label{eq:axm-local-feedback-links}
\end{equation}
Whenever \(s_M<1\), human research weakens the local capability feedback. If
AI research services eventually dominate research expenditure, then
\(s_M\to1\) and the extension approaches the benchmark feedback.

Given \(K,A,B,N,q\), the extension's algebraic subsystem consists of labor
clearing; final production and prices; the service identity; inverse demand and
the monopoly condition; and the research price, conditional demands, and scale
condition. On the selected interior branch, it determines allocations
\(L,H,U,M,X,E,Z\), output \(Y\), prices \(r,w,p_X,P_E\), and shares \(s_X,s_M\).
The dynamic system is
\begin{subequations}
\label{eq:axm-canonical-system}
\begin{align}
 \frac{\dot N}{N}&=n,\\
 \frac{\dot A}{A}&=\gamma_A,\\
 \frac{\dot K}{K}&=\frac{Y-C-U-M}{K}-\delta,\\
 \frac{\dot B}{B}&=\frac{\chi E^\eta}{B},\\
 \frac{\dot C}{C}&=n+\alpha\frac{Y}{K}-\delta-\rho,\\
 \frac{\dot q}{q}
 &=\alpha\frac{Y}{K}-\delta-\frac{X}{qB^2}
 -\eta s_M\frac{\dot B}{B}.
\end{align}
\end{subequations}

\begin{definition}[Interior equilibrium with human research]
Given \(K(0),A(0),B(0),N(0)>0\), an interior equilibrium with human research is
an admissible path with positive states and controls; in particular,
\(C,K,A,B,N\) and \(L,H,U,M,X,E,Z,Y\) are positive, as are
\(w,p_X,P_E,q\). The household, final-good firms, and developer solve their
respective problems; all markets clear; and the two transversality conditions in
\eqref{eq:axm-tvc} hold.
\end{definition}

The dated first-order, state, market, and accounting conditions are necessary
for a regular interior equilibrium. I call a path satisfying those conditions
and the two terminal restrictions a candidate path. The monopoly second-order
condition establishes only local optimality in \(X\), and I do not establish
global concavity of the developer's problem for all parameter values.

\subsection{Research curvature and sufficiency}

The restriction \(\eta<1\) gives decreasing returns to current research inputs
when capability is fixed. It need not bound the developer's value because
research also raises capability and therefore future operating profit. The next
result states the relevant large-scale restriction; the Appendix provides the
exact bounds and proof.

\begin{proposition}[Large-scale research expenditure]
\label{prop:axm-research-burst}
Suppose \(\sigma_{XL}>0\) and \(0<\eta<1\). Fix a finite horizon, an initial
capability stock, positive continuous paths for \(K,A,L\), and an integrable
path for \(r\). In the benchmark, consider all nonnegative measurable research-
compute policies with finite expenditure. If \(\eta<\alpha\), the truncated
developer value is finite and the objective tends to minus infinity as total
research expenditure diverges. If \(\sigma_{XL}>1\) and \(\eta>\alpha\), fixed-
duration bursts of research compute make the truncated value unbounded above.

The same finite-value conclusion holds in the human-research extension for any
\(\sigma_{HM}>0\), taking a positive continuous wage path as given. The
unbounded-value conclusion holds for the extension when
\(\sigma_{XL}>1\) and \(\sigma_{HM}>1\). At \(\eta=\alpha\), the leading benefit
and cost terms have the same order, so their coefficients determine the result
for this family of policies.
\end{proposition}

\begin{assumption}[Baseline large-scale research curvature]
\label{ass:axm-large-scale-curvature}
All equilibrium results and numerical exercises maintain
\begin{equation}
 0<\eta<\alpha<1.
 \label{eq:axm-large-scale-curvature}
\end{equation}
\end{assumption}

The restriction is stronger than decreasing returns to current research scale.
It rules out arbitrarily large research expenditure as optimal on every fixed
horizon satisfying the proposition's conditions. It does not prove that a
maximizer exists, establish global concavity, control the infinite-horizon tail,
or prove equilibrium existence. The numerical exercises use
\(\eta=0.20<\alpha=0.33\), chosen illustratively rather than estimated.

\begin{proposition}[Global sufficiency in the unit-elasticity cases]
\label{prop:axm-equilibrium-sufficiency}
Let \(\sigma_{XL}=1\) and
\(\beta=(1-\alpha)\omega_X\leq1/2\). Suppose one of the following conditions
holds:
\begin{enumerate}
 \item the automated-research benchmark applies and \(2\eta\leq1\);
 \item the human-research extension has
 \(\sigma_{HM}=1\) and \(\eta(1+\omega_M)\leq1\); or
 \item the human-research extension has
 \(\sigma_{HM}=2\) and \(2\eta\leq1\).
\end{enumerate}
Every regular equilibrium that satisfies the two transversality conditions
obeys the corresponding candidate system. Conversely, every admissible
infinite-horizon candidate path is an equilibrium.
\end{proposition}

The calibration satisfies these restrictions with \(\beta=0.134\),
\(2\eta=0.40\), and, in the unit-elasticity human-research extension,
\(\eta(1+\omega_M)=0.2700\). The proposition does not turn a finite-horizon
approximation into a proof that an admissible infinite-horizon continuation
exists or satisfies the two transversality restrictions.
